\newpage
\section*{Willkommen zur FOSSGIS-Konferenz 2026 in Göttingen!}\label{welcome}

\begin{flushright}
\includegraphics[width=1.0\textwidth]{Logo_Verein_25.png}
\end{flushright}

\noindent
Im Jahr 2025 wurde der FOSSGIS e.V. 25 Jahre alt. Gegründet wurde der Verein als
GRASS Anwendervereinung (GAV e.V.) am 15.04.2000 in Hannover. Über die Zeit
kamen immer mehr Software-Projekte hinzu, weshalb der Verein im Jahr 2008 in
FOSSGIS e.V. umbenannt wurde. {\bfseries  FOSSGIS} steht für {\bfseries f}reie
{\bfseries O}pen"={\bfseries S}ource"={\bfseries S}oftware für
{\bfseries G}eo\-{\bfseries i}nformations\-{\bfseries s}ysteme.
Neben dem Thema \glqq Open-Source-Software\grqq wurde auch das
Thema der offenen Geodaten immer wichtiger. Seit 2017 ist der FOSSGIS~e.V.
auch der offizielle Vertreter des OpenStreetMap-Projektes in Deutschland.

\noindent
Kern unserer Arbeit ist die jährliche FOSSGIS-Konferenz, die Entwickler und
Nutzende zusammenbringt. Wir sind aber auch über die Konferenzen hinaus aktiv,
gehen zu Veranstaltungen, beantworten Fragen in Online-Foren oder betreuen
die digitale Infrastruktur hinter den Projekten.

\pagebreak
\noindent
Die Nachfrage nach Informationen und Vernetzung steigt stetig an. Dies zeigt sich
an den Teilnehmerzahlen der Konferenzen, wie auch an der Zahl der E-Mails und
Telefonanrufe, die uns täglich erreichen. Daher hat der Verein seit fünf Jahren
eine professionelle Koordinierungsstelle, die die ehrenamtliche Arbeit
unterstützt und als Schnittstelle nach Außen dient. Ende 2023 wurde die
Koordinierungsstelle speziell für den OSM-Bereich erweitert.

\noindent
Nach wie vor wird die meiste Arbeit im Verein von Ehrenamtlichen getragen.
Dazu braucht es Menschen, die mitmachen und sich engagieren, Ideen einbringen
und unsere Zukunft mitgestalten wollen. Vielleicht kommen Sie mal am
FOSSGIS-Stand vorbei und reden Sie mit uns! Wir freuen uns auf Ihre Fragen und
Ideen und vielleicht haben Sie ja auch Lust mitzumachen!

\section*{Fotos und Social Media}
Wir freuen uns, wenn Sie Beiträge über die FOSSGIS-Konferenz in Social Media posten! Bitte verwenden Sie als Hashtags {\em \#fossgis2026 und ggf. \#openstreetmap}. Falls Sie Fotos von Personen veröffentlichen denken Sie bitte daran, vorher deren Einverständnis einzuholen.
\newpage
\section*{Exkursionen {\normalfont\em (Anmeldung erforderlich)}}\label{exkursionen}
\noindent
{\large \bfseries Führung durch das Rechenzentrum der GWDG}\\
Das Rechenzentrum der Gesellschaft für wissenschaftliche Datenverarbeitung mbH Göttingen (GWDG) fungiert als Hochschulrechenzentrum für die Georg-August-Universität Göttingen und Rechen- und IT-Kompetenzzentrum für die Max-Planck-Gesellschaft. Die Führung gibt einen kurzen Überblick die Dienstleistungen der GWDG sowie über die technische Ausstattung.
 {\bfseries Mittwoch, 25.03.2025, 17:00 und 13:15} Treffpunkt: Rechenzentrum der GWDG (RZGö) (vom ZHG mittels der Buslinien 21,22,23 innerhalb von ca. 10~Minuten erreichbar).
\bigskip

\noindent
{\large \bfseries Exkursion -- Führung durch das Forum Wissen}\\
Führung durch das Wissensmuseum der Universität Göttingen. Das Forum Wissen ist ein Ort der Begegnung und es geht um die Frage, wie Wissen entsteht.
{\bfseries Mittwoch, 25.03.2026 und Donnerstag, 26.03.2026, 15:00}, Treffpunkt: Forum Wissen, Berliner Straße~28 (nahe Hauptbahnhof)
\newpage
\noindent
{\large \bfseries Archäologische Stadtführung -- Schwerpunkt Stadtbefestigung}\\
Die Gruppe wird zweigeteilt, die eine Hälfte wird in der Düsteren Straße 21 archäologische Funde aus Göttingen kennenlernen und die andere einen kurzen Rundgang durch die Stadt machen und die Befestigungsanlagen und Fundstellen kennenlernen. Nach 45 Minuten werden die beiden Gruppen tauschen. {\bfseries Mittwoch, 25.03.2026, 15:30,
Donnerstag, 26.03.2026, 15:30}, Treffpunkt: Düstere Straße~21.
\bigskip

\noindent
{\large \bfseries Archäologische Stadtführung -- Schwerpunkt vom Dorf zur Stadt}\\
In einem eineinhalbstündigen Rundgang durch die Stadt werden die wichtigen Orte begangen, die Göttingen vom Dorf zur Stadt werden lassen. Angefangen bei dem \glqq Alten Dorf\grqq\  und der Albani Kirche über die spätere Stadtgründung mit dem Markt und den drei \glqq neuen\grqq\  Pfarrbezirken in der heutigen Stadtmitte. {\bfseries Freitag, 27.03.2026, 15:30}, Treffpunkt: Düstere Straße~21.
\bigskip

\newpage
\noindent
{\large \bfseries Exkursion -- Besondere Stücke der Kartensammlung der SUB Göttingen}\\
Seit Gründung der Niedersächsischen Staats- und Universitätsbibliothek Göttingen sind Karten ein besonders gepflegter Sammelschwerpunkt. Neben laufenden Einzelkäufen war es der Bibliothek immer wieder möglich, durch großzügige Schenkungen kostbare Sammlungen zu erwerben. In der Führung werden besondere Stücke gezeigt. {\bfseries Freitag, 27.03.2026, 17:00}, Treffpunkt: Historisches Gebäude, Papendiek~14, Haupteingang.
\bigskip

\noindent
{\large \bfseries Exkursion - Physicalisches Cabinett}\\
Die Sammlung historischer physikalischer Instrumente am I.~Physikalischen Institut, kurz Physicalisches Cabinet, erlaubt ihren Besuchern anhand von historischen Messapparaturen einen Einblick in die Geschichte der Physik an der Georg-August-Universität Göttingen.
{\bfseries Samstag, 28.03.2026, 10:30}, Treffpunkt: Eingang I.~Physikalisches Institut,
Friedrich-Hund-Platz~1 (gegenüber Geographischen Institut)

\newpage
\section*{Abendveranstaltung am Mittwoch}\label{schwaetzli}
Am ersten Abend der FOSSGIS-Konferenz findet die Abendveranstaltung im 1.~OG des  Foyer des Zentralen Hörsaalgebäudes von {\bfseries 19:00 bis 22:00 Uhr} statt. Eine Anmeldung ist erforderlich.

\section*{Rahmenprogramm am Donnerstag}
\subsection*{Gruppenfoto}
Wir laden am {\bfseries Donnerstag} in der {\bfseries Nachmittagspause} zum Gruppenfoto auf dem Platz vor dem Zentralen Hörsaalgebäude ein.

\subsection*{Mitgliederversammlung des FOSSGIS e.V.}
Am {\bfseries Donnerstag} sind alle Mitglieder und Gäste ab {\bfseries 19.00 Uhr} herzlich zur Mitgliederversammlung des FOSSGIS e.V. eingeladen.  Es gibt Getränke und Pizza. Der Verein freut sich über zahlreiches Erscheinen.

\newpage
\section*{Rahmenprogramm am Freitag}

\subsection*{Sektempfang am FOSSGIS-Stand}
Alle Mitglieder des FOSSGIS-Vereins, Freunde und Interessierte sind am {\bfseries Freitag} ab {\bfseries 16.15 Uhr} herzlich zum Sektempfang zum Ausklang der FOSSGIS 2026 am FOSSGIS-Vereins-Stand eingeladen.

\subsection*{OSM-Event am Freitagabend}
Zur Einstimmung auf den OSM-Samstag sind am {\bfseries Freitagabend}  ab 18:30 Uhr im {\em Kartoffelhaus} Plätze reserviert. Adresse: Goethe\-allee~8. Es gibt Essen sowie Getränke, jede:r bezahlt seine Rechnung selbst. Zur besseren Planung wird um Eintragung im OSM-Wiki gebeten.
\newpage
\section*{Programm am Samstag}

\subsection*{OSM-Samstag}
Am Samstag, den 28.~März findet von 9.00 bis 17.30~Uhr die OSM-Unkonferenz statt.
Interessierte sind eingeladen sich zu beteiligen oder daran teilzunehmen.
Die Themensammlung erfolgt im OSM-Wiki. Die Veranstaltung ist kostenfrei.
Um Anmeldung über das FOSSGIS-Konferenz-Anmeldesystem wird gebeten.

\subsection*{Community-Sprint}
Beim Community-Sprint wird gemeinsam an OpenSource Projekten gearbeitet.
Die Veranstaltung startet um 9~Uhr mit einem kurzen Vortrag, der für Einsteiger:innen
und Interessierte erklärt wie OpenSource funktioniert und wie man beitragen kann.
Im Anschluss wird gemeinsam oder individuell an Projekten gearbeitet. Themen können im OSM-Wiki
eingetragen werden. Es ist auch eine Beteiligung an der Initiative "QGIS-Plugins für Qt6/QGIS4" möglich.
Um Anmeldung über das FOSSGIS-Konferenz-Anmeldesystem wird gebeten.
